\documentclass[12pt]{article}
% Page setup
\usepackage[margin=1in]{geometry}
\usepackage{parskip}
\usepackage{titling}
\setlength{\droptitle}{-6em} 

% Math packages
\usepackage{amsmath, amsthm, amssymb, amsfonts}
\usepackage{mathtools}
\makeatletter
\g@addto@macro\th@plain{\normalfont} % remove italics from theorems (plain style)
\makeatother

% Commutative diagrams
\usepackage{tikz-cd}


% Theorem environments
\newtheorem{theorem}{Theorem}[subsection]
\newtheorem{lemma}[theorem]{Lemma}
\newtheorem{proposition}[theorem]{Proposition}
\newtheorem{corollary}[theorem]{Corollary}
\newtheorem{fact}[theorem]{Fact}
\newtheorem{remark}[theorem]{Remark}

\theoremstyle{definition}
\newtheorem{definition}[theorem]{Definition}
\newtheorem{example}[theorem]{Example}

% Useful commands for category theory
\newcommand{\mc}{\mathcal}
\newcommand{\ini}{\operatorname{ini}}
\newcommand{\ter}{\operatorname{ter}}
\newcommand{\cod}{\operatorname{cod}}
\newcommand{\dom}{\operatorname{dom}}

\newcommand{\List}{\operatorname{List}}
\newcommand{\lop}{\operatorname{[ \ ]}}
\newcommand{\len}{\operatorname{len}}
\newcommand{\type}{\operatorname{type}}

\newcommand{\Exc}{\operatorname{Exception}}

\newcommand{\Kl}{\operatorname{Kl}}
\newcommand{\plus}{\operatorname{+}}

\newcommand{\Cat}{\mathcal{C}at}
\newcommand{\Vect}{\mathrm{Vect}_\mathbf{k}}
\newcommand{\Hom}{\operatorname{Hom}}
\newcommand{\Nat}{\operatorname{Nat}}
\newcommand{\id}{\mathrm{id}}
\newcommand{\op}{\mathrm{op}}
\newcommand{\obj}{\operatorname{obj}}
\newcommand{\arr}{\operatorname{arr}}
\newcommand{\adj}{\operatorname{adj}}
\newcommand{\coadj}{\operatorname{coadj}}
\newcommand{\colim}{\operatorname*{colim}}

% Arrows
\newcommand{\ra}{\rightarrow}
\newcommand{\from}{\leftarrow}
\newcommand{\To}{\Rightarrow}
\newcommand{\xto}{\xrightarrow}
\newcommand{\xfrom}{\xleftarrow}

\renewcommand{\labelitemi}{--} % Level 1: Dashes
\renewcommand{\labelitemii}{$\circ$} % Level 2: Circles
\renewcommand{\labelitemiii}{$\bullet$} % Level 3: Bullets (default)


\title{Monoidal categories}
\author{Alyson Mei}
\date{\today}

\begin{document}


\maketitle

\tableofcontents

\paragraph{Conventions}
\hfill
\begin{itemize}
    \item We work in the arrow-first categorical setting;
    \item Lambda-abstraction and left unitor use the same letter $\lambda$, but the first one is used only in the form of $\lambda x.y$, so the local meaning is clear;
    \item Functors of the form $\lambda f_1...f_n.F$ are typed $\mc{C}^n \to \mc{D}$ with coherence for product's associativity considered to be known.
\end{itemize}

\newpage

\section{Core definitions and results}

\subsection{Core definitions}

\subsubsection{Monoidal categories: basic types}

\begin{definition}[Monoidal category]
    \emph{Monoidal category (MC)} is a tuple
    $$(\mc{C}, \otimes, I, \alpha, \lambda, \rho),$$
    where 
    \begin{itemize}
        \item[$\bullet$] $\mc{C}$ is a \emph{base category},
        \item[$\bullet$] $\otimes: \mc{C} \times \mc{C} \to \mc{C}$ is a \emph{monoidal product} (or \emph{tensor}), 
        \item[$\bullet$] $I : \obj(\mc{C})$ is a \emph{monoidal unit},
        \item[$\bullet$] $\alpha: \lambda fgh. (f \otimes g) \otimes h \xto{\simeq} \lambda fgh.f \otimes (g \otimes h)$ is a \emph{associator},
        \item[$\bullet$] $\lambda: \lambda f. I \otimes f \xto{\simeq} \id_{\mc{C}}$ is a \emph{left unitor},
        \item[$\bullet$] $\rho: \lambda f. f \otimes I \xto{\simeq} \id_{\mc{C}}$ is a \emph{right unitor}, 
    \end{itemize}
    satisfying the \emph{associativity (pentagon)} and \emph{unit (triangle)} diagrams:
    % https://q.uiver.app/#q=WzAsNSxbMCwwLCJhIFxcb3RpbWVzIChiIFxcb3RpbWVzIChjIFxcb3RpbWVzIGQpKSJdLFsxLDAsIihhIFxcb3RpbWVzIGIpIFxcb3RpbWVzIChjIFxcb3RpbWVzIGQpIl0sWzIsMCwiKChhIFxcb3RpbWVzIGIpIFxcb3RpbWVzIGMpIFxcb3RpbWVzIGQiXSxbMCwxLCJhIFxcb3RpbWVzICgoYiBcXG90aW1lcyBjKSBcXG90aW1lcyBkKSJdLFsyLDEsIihhIFxcb3RpbWVzIChiIFxcb3RpbWVzIGMpKSBcXG90aW1lcyBkIl0sWzAsMSwiXFxhbHBoYV97YSwgYiwgKGMgXFxvdGltZXMgZCl9Il0sWzEsMiwiXFxhbHBoYV97KGEgXFxvdGltZXMgYiksIGMsIGR9Il0sWzAsMywiXFxpZF9hIFxcb3RpbWVzIFxcYWxwaGFfe2IsIGMsIGR9IiwyXSxbMyw0LCJcXGFscGhhX3thLCBiXFxvdGltZXMgYywgZH0iLDJdLFs0LDIsIlxcYWxwaGFfe2EsIGIsIGN9IFxcb3RpbWVzIFxcaWRfZCIsMl1d
    \[\begin{tikzcd}
        {a \otimes (b \otimes (c \otimes d))} & {(a \otimes b) \otimes (c \otimes d)} & {((a \otimes b) \otimes c) \otimes d} \\
        {a \otimes ((b \otimes c) \otimes d)} && {(a \otimes (b \otimes c)) \otimes d}
        \arrow["{\alpha_{a, b, (c \otimes d)}}", from=1-1, to=1-2]
        \arrow["{\id_a \otimes \alpha_{b, c, d}}"', from=1-1, to=2-1]
        \arrow["{\alpha_{(a \otimes b), c, d}}", from=1-2, to=1-3]
        \arrow["{\alpha_{a, b\otimes c, d}}"', from=2-1, to=2-3]
        \arrow["{\alpha_{a, b, c} \otimes \id_d}"', from=2-3, to=1-3]
    \end{tikzcd},\]
    % https://q.uiver.app/#q=WzAsMyxbMCwwLCJhIFxcb3RpbWVzIChJIFxcb3RpbWVzIGIpIl0sWzIsMCwiKGEgXFxvdGltZXMgSSkgXFxvdGltZXMgYiJdLFsxLDEsImEgXFxvdGltZXMgYiJdLFswLDEsIlxcYWxwaGFfe2EsIEksIGJ9Il0sWzAsMiwiXFxpZF9hIFxcb3RpbWVzIFxcbGFtYmRhX2IiLDJdLFsyLDEsIlxccmhvX2EgXFxvdGltZXMgXFxpZF9iIiwyXV0=
    \[\begin{tikzcd}
        {a \otimes (I \otimes b)} && {(a \otimes I) \otimes b} \\
        & {a \otimes b}
        \arrow["{\alpha_{a, I, b}}", from=1-1, to=1-3]
        \arrow["{\id_a \otimes \lambda_b}"', from=1-1, to=2-2]
        \arrow["{\rho_a \otimes \id_b}"', from=2-2, to=1-3]
    \end{tikzcd}\text{.}\]
\end{definition}

\begin{definition}[Strict monoidal category]
    \emph{Strict monoidal category} is an MC with $\alpha$, $\rho$ and $\lambda$ being identities. Thus,
    \begin{itemize}
        \item $(a \otimes b) \otimes c = a \otimes (b \otimes c)$,
        \item $I \otimes a = a \otimes I = a.$
    \end{itemize}
\end{definition}

\begin{definition}[Braided monoidal category]
\label{def:braided_mc}
    \emph{Strict monoidal category} is an MC equipped with \emph{braiding} $$\sigma: \lambda fg. f \otimes g \to \lambda fg. g \otimes f$$ satisfying the \emph{hexagon} diagrams:
    % https://q.uiver.app/#q=WzAsNixbMSwwLCJhIFxcb3RpbWVzIChiIFxcb3RpbWVzIGMpIl0sWzIsMCwiKGIgXFxvdGltZXMgYykgXFxvdGltZXMgYSJdLFszLDEsImIgXFxvdGltZXMgKGMgXFxvdGltZXMgYSkiXSxbMCwxLCIoYSBcXG90aW1lcyBiKSBcXG90aW1lcyBjIl0sWzEsMiwiKGIgXFxvdGltZXMgYSkgXFxvdGltZXMgYyJdLFsyLDIsImIgXFxvdGltZXMgKGEgXFxvdGltZXMgYykiXSxbMSwyLCJcXGFscGhhIl0sWzAsMSwiXFxnYW1tYSJdLFszLDAsIlxcYWxwaGEiXSxbMyw0LCJcXGdhbW1hIFxcb3RpbWVzIFxcaWQiLDJdLFs0LDUsIlxcYWxwaGEiLDJdLFs1LDIsIlxcaWQgXFxvdGltZXMgXFxnYW1tYSIsMl1d
    \[\begin{tikzcd}[sep=scriptsize]
        & {a \otimes (b \otimes c)} & {(b \otimes c) \otimes a} \\
        {(a \otimes b) \otimes c} &&& {b \otimes (c \otimes a)} \\
        & {(b \otimes a) \otimes c} & {b \otimes (a \otimes c)}
        \arrow["\sigma", from=1-2, to=1-3]
        \arrow["\alpha", from=1-3, to=2-4]
        \arrow["\alpha", from=2-1, to=1-2]
        \arrow["{\sigma \otimes \id}"', from=2-1, to=3-2]
        \arrow["\alpha"', from=3-2, to=3-3]
        \arrow["{\id \otimes \sigma}"', from=3-3, to=2-4]
    \end{tikzcd},\]
    % https://q.uiver.app/#q=WzAsNixbMSwwLCIoYSBcXG90aW1lcyBiKSBcXG90aW1lcyBjIl0sWzIsMCwiYyBcXG90aW1lcyAoYSBcXG90aW1lcyBiKSJdLFszLDEsIihjIFxcb3RpbWVzIGEpIFxcb3RpbWVzIGIiXSxbMCwxLCJhIFxcb3RpbWVzIChiIFxcb3RpbWVzIGMpIl0sWzEsMiwiYSBcXG90aW1lcyAoYyBcXG90aW1lcyBiKSJdLFsyLDIsIihhIFxcb3RpbWVzIGMpIFxcb3RpbWVzIGIiXSxbMSwyLCJcXGFscGhhXnstMX0iXSxbMCwxLCJcXGdhbW1hIl0sWzMsMCwiXFxhbHBoYV57LTF9Il0sWzMsNCwiXFxpZCBcXG90aW1lcyBcXGdhbW1hIiwyXSxbNCw1LCJcXGFscGhhXnstMX0iLDJdLFs1LDIsIlxcZ2FtbWEgXFxvdGltZXMgXFxpZCIsMl1d
    \[\begin{tikzcd}[sep=scriptsize]
        & {(a \otimes b) \otimes c} & {c \otimes (a \otimes b)} \\
        {a \otimes (b \otimes c)} &&& {(c \otimes a) \otimes b} \\
        & {a \otimes (c \otimes b)} & {(a \otimes c) \otimes b}
        \arrow["\sigma", from=1-2, to=1-3]
        \arrow["{\alpha^{-1}}", from=1-3, to=2-4]
        \arrow["{\alpha^{-1}}", from=2-1, to=1-2]
        \arrow["{\id \otimes \sigma}"', from=2-1, to=3-2]
        \arrow["{\alpha^{-1}}"', from=3-2, to=3-3]
        \arrow["{\sigma \otimes \id}"', from=3-3, to=2-4]
    \end{tikzcd}
    \text{.}
    \]
\end{definition}

\begin{definition}[Symmetric monoidal category]
    \hfill\break
    \emph{Symmetric monoidal category (SMC)} can be defined in two equivalent ways:
    \begin{enumerate}
        \item[a)] SMC is a braided MC with braiding $\sigma$ being iso and thus called a \emph{swap}:
        $$\sigma_{f, g}: f \otimes g \xto{\cong} g \otimes f;$$
        \item[b)] SMC is an MC with a \emph{swap} iso
        $$\sigma: \lambda fg. f \otimes g \xto{\simeq} \lambda fg. g \otimes f,$$
        satisfying the \emph{associativity coherence} (that is exactly the first hexagon diagram from the Def \ref{def:braided_mc}) with additional \emph{unit coherence} and \emph{inverse law} diagrams:
        % https://q.uiver.app/#q=WzAsNixbMCwwLCJhIFxcb3RpbWVzIEkiXSxbMiwwLCJJIFxcb3RpbWVzIGEiXSxbMSwxLCJhIl0sWzMsMSwiYSBcXG90aW1lcyBiIl0sWzUsMSwiYSBcXG90aW1lcyBiIl0sWzQsMCwiYiBcXG90aW1lcyBhIl0sWzAsMSwiXFxzaWdtYV97YSwgSX0iXSxbMCwyLCJcXHJob19hIiwyXSxbMSwyLCJcXGxhbWJkYV9iIl0sWzMsNCwiXFxpZF97YSBcXG90aW1lcyBifSIsMl0sWzMsNSwiXFxzaWdtYV97YSwgYn0iXSxbNSw0LCJcXHNpZ21hX3tiLCBhfSJdXQ==
        \[\begin{tikzcd}
            {a \otimes I} && {I \otimes a} && {b \otimes a} \\
            & a && {a \otimes b} && {a \otimes b}
            \arrow["{\sigma_{a, I}}", from=1-1, to=1-3]
            \arrow["{\rho_a}"', from=1-1, to=2-2]
            \arrow["{\lambda_b}", from=1-3, to=2-2]
            \arrow["{\sigma_{b, a}}", from=1-5, to=2-6]
            \arrow["{\sigma_{a, b}}", from=2-4, to=1-5]
            \arrow["{\id_{a \otimes b}}"', from=2-4, to=2-6]
        \end{tikzcd}
        \text{.}
        \]
    \end{enumerate}
\end{definition}

\subsubsection{Monoidal functors}

\begin{definition}[(Op)lax monoidal functor]
    Let $(\mc{C}, \otimes, I_{\mc{C}})$ and $(\mc{C}, \bullet, I_{\mc{D}})$ be monoidal.
    \begin{itemize}
        \item A \emph{lax monoidal functor} (or just \emph{monoidal functor}) from $\mc{C}$ to $\mc{D}$ is a triple $$(F, \phi, \phi_0),$$
        where \begin{itemize}
            \item[$\bullet$] $F: \mc{C} \to \mc{D}$,
            \item[$\bullet$] $\phi: \lambda fg. Ff \bullet Fg \to \lambda fg. F(f \otimes g), $
            \item[$\bullet$] $\phi_0: I_{\mc{D}} \to FI_{\mc{C}},$ 
        \end{itemize}
        satisfying the diagrams:
        % https://q.uiver.app/#q=WzAsNixbMCwxLCIoRmEgXFxidWxsZXQgRmIpIFxcYnVsbGV0IEZjIl0sWzEsMiwiRihhIFxcb3RpbWVzIGIpIFxcYnVsbGV0IEZjIl0sWzEsMCwiRmEgXFxidWxsZXQgKEZiIFxcYnVsbGV0IEZjKSJdLFsyLDIsIkYoKGEgXFxvdGltZXMgYikgXFxvdGltZXMgYykiXSxbMywxLCJGKGEgXFxvdGltZXMgKGIgXFxvdGltZXMgYykpIl0sWzIsMCwiRmEgXFxidWxsZXQgRihiIFxcb3RpbWVzIGMpIl0sWzAsMSwiXFxwaGkgXFxidWxsZXQgXFxpZCIsMl0sWzAsMiwiXFxhbHBoYV97XFxtYXRoY2Fse0R9fSJdLFsxLDMsIlxccGhpIiwyXSxbMyw0LCJGXFxhbHBoYV97XFxtYXRoY2Fse0N9fSIsMl0sWzIsNSwiXFxpZCBcXGJ1bGxldCBcXHBoaSJdLFs1LDQsIlxccGhpIl1d
        \[\begin{tikzcd}[sep=scriptsize]
            & {Fa \bullet (Fb \bullet Fc)} & {Fa \bullet F(b \otimes c)} \\
            {(Fa \bullet Fb) \bullet Fc} &&& {F(a \otimes (b \otimes c))} \\
            & {F(a \otimes b) \bullet Fc} & {F((a \otimes b) \otimes c)}
            \arrow["{\id \bullet \phi}", from=1-2, to=1-3]
            \arrow["\phi", from=1-3, to=2-4]
            \arrow["{\alpha_{\mathcal{D}}}", from=2-1, to=1-2]
            \arrow["{\phi \bullet \id}"', from=2-1, to=3-2]
            \arrow["\phi"', from=3-2, to=3-3]
            \arrow["{F\alpha_{\mathcal{C}}}"', from=3-3, to=2-4]
        \end{tikzcd},\]
        % https://q.uiver.app/#q=WzAsOCxbMywxLCJGYSBcXGJ1bGxldCBJX3tcXG1hdGhjYWx7RH19Il0sWzQsMSwiRmEgXFxidWxsZXQgRklfe1xcbWF0aGNhbHtEfX0iXSxbMywwLCJGYSJdLFs0LDAsIkYoYSBcXG90aW1lcyBJX3tcXG1hdGhjYWx7Q319KSJdLFsxLDAsIklfXFxtYXRoY2Fse0R9IFxcYnVsbGV0IEZiIl0sWzAsMCwiRklfe1xcbWF0aGNhbHtDfX0gXFxidWxsZXQgRmIiXSxbMSwxLCJGYiJdLFswLDEsIkYoSV9cXG1hdGhjYWx7Q30gXFxvdGltZXMgQikiXSxbMCwyLCJcXHJob197XFxtYXRoY2Fse0R9fSJdLFszLDIsIkZcXHJob197XFxtYXRoY2Fse0N9fSIsMl0sWzEsMywiXFxwaGkiLDJdLFswLDEsIlxcaWQgXFxidWxsZXQgXFxwaGkiLDJdLFs0LDYsIlxcbGFtYmRhX3tcXG1hdGhjYWx7RH19Il0sWzQsNSwiXFxwaGkgXFxidWxsZXQgXFxpZCIsMl0sWzUsNywiXFxwaGkiLDJdLFs3LDYsIkZcXGxhbWJkYV97XFxtYXRoY2Fse0N9fSIsMl1d
    % https://q.uiver.app/#q=WzAsOCxbMywxLCJGYSBcXGJ1bGxldCBJX3tcXG1hdGhjYWx7RH19Il0sWzQsMSwiRmEgXFxidWxsZXQgRklfe1xcbWF0aGNhbHtEfX0iXSxbMywwLCJGYSJdLFs0LDAsIkYoYSBcXG90aW1lcyBJX3tcXG1hdGhjYWx7Q319KSJdLFsxLDAsIklfXFxtYXRoY2Fse0R9IFxcYnVsbGV0IEZiIl0sWzAsMCwiRklfe1xcbWF0aGNhbHtDfX0gXFxidWxsZXQgRmIiXSxbMSwxLCJGYiJdLFswLDEsIkYoSV9cXG1hdGhjYWx7Q30gXFxvdGltZXMgQikiXSxbMCwyLCJcXHJob197XFxtYXRoY2Fse0R9fSJdLFszLDIsIkZcXHJob197XFxtYXRoY2Fse0N9fSIsMl0sWzEsMywiXFxwaGkiLDJdLFswLDEsIlxcaWQgXFxidWxsZXQgXFxwaGlfMCIsMl0sWzQsNiwiXFxsYW1iZGFfe1xcbWF0aGNhbHtEfX0iXSxbNCw1LCJcXHBoaV8wIFxcYnVsbGV0IFxcaWQiLDJdLFs1LDcsIlxccGhpIiwyXSxbNyw2LCJGXFxsYW1iZGFfe1xcbWF0aGNhbHtDfX0iLDJdXQ==
    \[\begin{tikzcd}
        {FI_{\mathcal{C}} \bullet Fb} & {I_\mathcal{D} \bullet Fb} && Fa & {F(a \otimes I_{\mathcal{C}})} \\
        {F(I_\mathcal{C} \otimes B)} & Fb && {Fa \bullet I_{\mathcal{D}}} & {Fa \bullet FI_{\mathcal{D}}}
        \arrow["\phi"', from=1-1, to=2-1]
        \arrow["{\phi_0 \bullet \id}"', from=1-2, to=1-1]
        \arrow["{\lambda_{\mathcal{D}}}", from=1-2, to=2-2]
        \arrow["{F\rho_{\mathcal{C}}}"', from=1-5, to=1-4]
        \arrow["{F\lambda_{\mathcal{C}}}"', from=2-1, to=2-2]
        \arrow["{\rho_{\mathcal{D}}}", from=2-4, to=1-4]
        \arrow["{\id \bullet \phi_0}"', from=2-4, to=2-5]
        \arrow["\phi"', from=2-5, to=1-5]
    \end{tikzcd}\text{.}\]
    \item An \emph{oplax monoidal functor} (or \emph{comonoidal functor}) from $\mc{D}$ to $\mc{C}$ is a dual to lax monoidal functor.
    \end{itemize}
\end{definition}

\begin{definition}[Strong monoidal functor]
    A \emph{strong monoidal functor} is a monoidal functors with the coherence maps $\phi$ and $\phi_0$ being iso. Thus,
    $$Ff \otimes Fg \cong F(f \otimes g), \qquad I_{\mc{D}} \cong FI_{\mc{C}}.$$ 
\end{definition}

\begin{definition}[Strict monoidal functor]
    A \emph{strict monoidal functor} is a monoidal functor with the coherence maps $\phi$ and $\phi_0$ being identities. Thus,
    $$Ff \otimes Fg = F(f \otimes g), \qquad I_{\mc{D}} = FI_{\mc{C}}.$$ 
\end{definition}

\begin{definition}[Braided monoidal functor]
    A \emph{braided monoidal functor} is a monoidal functor between braided categories satisfying the diagram
    % https://q.uiver.app/#q=WzAsNCxbMSwwLCJGYSBcXGJ1bGxldCBGYiJdLFswLDAsIkZhIFxcYnVsbGV0IEZiIl0sWzAsMSwiRihhIFxcb3RpbWVzIGIpIl0sWzEsMSwiRihiIFxcb3RpbWVzIGEpIl0sWzEsMiwiXFxwaGkiLDJdLFsyLDMsIkZcXHNpZ21hX3tcXG1hdGhjYWx7RH19IiwyXSxbMSwwLCJcXHNpZ21hX3tcXG1hdGhjYWx7Q319Il0sWzAsMywiXFxwaGkiXV0=
    \[\begin{tikzcd}
        {Fa \bullet Fb} & {Fa \bullet Fb} \\
        {F(a \otimes b)} & {F(b \otimes a)}
        \arrow["{\sigma_{\mathcal{C}}}", from=1-1, to=1-2]
        \arrow["\phi"', from=1-1, to=2-1]
        \arrow["\phi", from=1-2, to=2-2]
        \arrow["{F\sigma_{\mathcal{D}}}"', from=2-1, to=2-2]
    \end{tikzcd}
    \text{.}
    \]
\end{definition}

\begin{definition}[Symmetric monoidal functor]
    A \emph{symmetric monoidal functor} is a braided monoidal functor whose domains are SMCs.
\end{definition}

\begin{definition}[Monoidal natural transformaiton]
    
\end{definition}

\subsection{Core results}


\section{Specialized definitions and results}


\end{document}
